\input{preamble.tex}
\usepackage{makecell} % commande \thead, dans l'exo 1
\usepackage{eurosym} % pour le symbole euro propre

\begin{document}

\reversemarginpar

\pagestyle{fancy}
\fancyhead[L]{Terminale STMG}
\fancyhead[C]{\textbf{DS n°4 - Fonction inverse et polynômes}}
\fancyhead[R]{\today}

\null\vspace{-30pt}
Nom / Prénom : \\

Consignes particulières : 
\begin{itemize}[label=$\bullet$]
	\item 
	La calculatrice est {autorisée}.
	\item 
	L'exercice \ref{exe:1a} peut être fait entièrement sur la feuille d'évaluation. Écrire son nom avant de rendre le sujet pour qu'il soit corrigé.
	\item 
	Toute trace de recherche est prise en compte.
%	\item 
%	L'abbréviation $\tq$ signifie ``tel que''.
\end{itemize}

\hrule

\exe{4}{
Compléter les phrases ci-dessous :
\begin{enumerate}
\item
La dérivée du polynôme $x \mapsto ax^3 + cx + d$ est $\dots$  \\
\item
La dérivée de la fonction $x \mapsto \dfrac{a}{x}$ sur $\Ret$ est $\dots$ \\
\item 
Le polynôme $8x - 3x^2 + 7$ est de degré $\dots$ \\
\item
La fonction $x \mapsto \dfrac{1}{x}$ est appelée $\dots$
\end{enumerate}
}{exe:1a}{

}

\exe{4}{
Donner les fonctions dérivées des fonctions suivantes sur $\Ret$ :
\begin{multicols}{2}
\begin{enumerate}
\item $f: x \mapsto 4x^2 - \frac5{x} + 3x$
\item $f: x \mapsto -12x^3  +\frac57 x^2 - 4 + \frac8{x}$
\item $f: x \mapsto x - x + x + \frac7{3x}$
\item$f: x \mapsto x^3 - \frac{3}2x^2 + 3x -3 - \frac{12}{x}$
\end{enumerate}
\end{multicols}
}{exe:2a}{

}

\exe{4}{
Etudier les variations des fonctions suivantes sur $\R$. Si possible, donner le maximum ou le minimum de la fonction.
\begin{enumerate}
\item $f: x \mapsto -3x^2 +18x + 32$
\item $f: x \mapsto \frac{15}2 x^2 + 5 x - 2$
\end{enumerate}
}{exe:3a}{
\begin{multicols}{2}
\begin{enumerate}
\item 
\item 
\item 
\item 
\end{enumerate}
\end{multicols}
}

\exe{8}{
Etudier les variations des fonctions suivantes sur $\Ret$ :
\begin{multicols}{2}
\begin{enumerate}
\item $f: x \mapsto 4x - 8 - \frac{12}x$
\item $f: x \mapsto 2x - 426 + \frac{242}{x}$
\item $f: x \mapsto-2x + 7 + \frac{5}{x}$
\item $f: x \mapsto-4x + 31 - \frac{144}{x}$
\end{enumerate}
\end{multicols}
}{exe:4a}{

}

%\newpage
%\fancyhead[C]{\textbf{Solutions}}
%Les solutions sont données dans l'ordre sujet a puis b. \newline
%\shipoutAnswer


\end{document}
